\documentclass[10pt,a4paper]{amsart}
\usepackage[margin=19mm]{geometry}
\usepackage{amsmath,amssymb,mathtools}
\usepackage[scr=rsfs,cal=euler]{mathalfa}
\usepackage{xfrac}
\usepackage[unicode,hidelinks,bookmarks=false]{hyperref}
\newcommand{\ie}{i.e.\ }
\newcommand{\cf}{cf.\ }
\newcommand{\resp}{respectively\ }
\newcommand{\Subsection}{\subsection}
\providecommand{\nobreakdash}{\nobreak\hbox{-}}
\setlength{\emergencystretch}{2em}
\multlinegap0pt
\usepackage{bm}
\usepackage{bbm}
\usepackage{dsfont}
\usepackage{xspace}
\usepackage{cancel}
\usepackage{mathtools}
\usepackage{amsthm}
\makeatletter
\@ifundefined{theorem}{\newtheorem{theorem}{Theorem}}{}
\makeatother
\title[Constant sign and nodal solutions for singular quasilinear e]{Constant sign and nodal solutions for singular quasilinear elliptic systems}
\author{Nouredine Medjoudj, Abdelkrim Moussaoui}
\date{Source: arXiv:2510.21489v1 (CC BY 4.0)}
\begin{document}
\maketitle

\noindent\emph{Source and licence.} Constant sign and nodal solutions for singular quasilinear elliptic systems. Version arXiv:2510.21489v1 (CC BY 4.0), distributed under the Creative Commons Attribution 4.0 International licence, as recorded in the arXiv licence metadata for this exact version and shown on the abstract page https://arxiv.org/abs/2510.21489v1. Full rights notice: licenses/arxiv-2510.21489-CC-BY-4.0.txt.\par\smallskip

\noindent\textit{Editorial scope.} Counted unit: the Theorem whose source label is \texttt{T4} in the article (source0.tex lines 1506--1516, source label \texttt{T4}), delivered together with the article's own complete original proof of it (source0.tex lines 1518--1543, one \texttt{proof} environment of 26 lines). No mathematics is written, repaired or re-derived: statement and proof are the article's own text line for line. Required context retained uncounted: T1 (lines 380--406); T2 (lines 649--659); 17 (lines 651--654). Every definition, display and label of those lines is retained unchanged. Omitted from this package and not redistributed: 1--379, 407--648, 655--1505, 1517--1517, 1544--1686. Nothing inside a retained excerpt is omitted or altered: every retained body is delivered exactly as the article's lines are written, so no omission receipt of any kind accompanies this package. Citations: the retained excerpts carry no citation command at all, so no bibliography entry of the article is transcribed and no reference is redistributed. Reference adaptation: none was needed and none was made; every reference inside the delivered unit resolves inside it, so no unresolved reference or citation is produced. Printed slips kept literal: no printed word, symbol, hypothesis or number of the counted statement or of its proof was altered. Layout and class: the article's own class is \texttt{amsart} with packages including eurosym, amssymb, amsfonts, amsmath, xcolor, bbm; the wrapper is the lane's pinned \texttt{amsart} editorial wrapper at 10pt on A4, which restates the theorem-like environments the retained text needs and adds the article's own macro definitions that the retained text needs (none), with extra wrapper packages (bm, bbm, dsfont, xspace, cancel, mathtools). The delivered numbering is the unit's own. Assets: no retained span contains a figure environment, a graphics-inclusion command, a PDF-page inclusion, a table or a TikZ picture; no image file is copied into this package, the package declares no asset dependency and no image file is redistributed with it. Licence: version arXiv:2510.21489v1 of this article is distributed under the Creative Commons Attribution 4.0 International licence, as recorded in the arXiv licence metadata for this exact version and shown on the abstract page https://arxiv.org/abs/2510.21489v1. A full rights notice is delivered at licenses/arxiv-2510.21489-CC-BY-4.0.txt. Characters: every retained excerpt is pure ASCII, so the delivered engine, XeLaTeX with its default Latin Modern text font, reports no missing character.
\begin{theorem}
\label{T1}Assume that $(\mathrm{H}_{2})$ and $(\mathrm{H}_{3})$ hold. Then,
problem $(\mathrm{P})$ admits two opposite constant-sign solutions $%
(u_{1,+},u_{2,+})$ and $(u_{1,-},u_{2,-})$ in $\mathcal{C}^{1,\tau }( 
\overline{\Omega })\times \mathcal{C}^{1,\tau }(\overline{\Omega }),$ for
certain $\tau \in (0,1).$ Moreover, every positive solution $%
(u_{1,+},u_{2,+})\in \mathcal{C}^{1}(\overline{\Omega })\times \mathcal{C}
^{1}(\overline{\Omega })$ and negative solution $(u_{1,-},u_{2,-})\in 
\mathcal{C}^{1}(\overline{\Omega })\times \mathcal{C}^{1}(\overline{\Omega }
) $ of $(\mathrm{P})$ within $[0,\overline{u}_{1}]\times \lbrack 0,\overline{
u}_{2}]$ and $[-\overline{u}_{1},0]\times \lbrack -\overline{u}_{2},0],$
respectively, satisfy 
\begin{equation}
u_{i,+}(x)\gg \underline{u}_{i}(x)\text{ \ and \ }u_{i,-}(x)\ll -\underline{%
u }_{i}(x),\text{ \ }\forall x\in \Omega .  \label{9}
\end{equation}
In particular, every positive solution $(u_{1,+},u_{2,+})\in
W_{0}^{1,p_{1}}(\Omega )\times W_{0}^{1,p_{2}}(\Omega )$ and negative
solution $(u_{1,-},u_{2,-})\in W_{0}^{1,p_{1}}(\Omega )\times
W_{0}^{1,p_{2}}(\Omega )$ of $(\mathrm{P})$ within $[0,\overline{u}
_{1}]\times \lbrack 0,\overline{u}_{2}]$ and $[-\overline{u}_{1},0]\times
\lbrack -\overline{u}_{2},0]$, respectively, satisfy 
\begin{equation}
u_{i,+}(x)\geq \underline{u}_{i}(x)\text{ \ and \ }u_{i,-}(x)\leq - 
\underline{u}_{i}(x),\text{ \ for a.e. }x\in \Omega .  \label{10}
\end{equation}
\end{theorem}

\begin{theorem}
\label{T2}Assume $(\mathrm{H}_{1})-(\mathrm{H}_{4})$ hold such that 
\begin{equation}
\alpha _{2}\geq \hat{\alpha}_{2}>0\text{ \ and }\ \beta _{1}\geq \hat{\beta}%
_{1}>0.  \label{17}
\end{equation}%
Then, system $(\mathrm{P})$ possesses nodal solutions $(u_{1}^{\ast
},u_{2}^{\ast })$ in $W_{0}^{1,p_{1}}(\Omega )\times W_{0}^{1,p_{2}}(\Omega
) $ where components $u_{1}^{\ast }$ and $u_{2}^{\ast }$ are nontrivial and
at least are not of the same constant sign.
\end{theorem}

\begin{theorem}
\label{T4} Assume $(\mathrm{H}_{1})-(\mathrm{H}_{5})$\textrm{\ }and (\ref{17}
) are fulfilled. Then, problem $(\mathrm{P})$ admits a nodal solution $%
(u_{1}^{\ast },u_{2}^{\ast })$ in $W_{0}^{1,p_{1}}(\Omega )\times
W_{0}^{1,p_{2}}(\Omega ),$ where both components $u_{1}^{\ast }$ and $%
u_{2}^{\ast }$ are sign changing. Moreover, $u_{1}^{\ast }$ and $u_{2}^{\ast
}$ change sign simultaneously, that is, 
\begin{equation}
u_{1}^{\ast }u_{2}^{\ast }\geq 0.  \label{55}
\end{equation}
\end{theorem}

\begin{proof}
Recall from Theorem \ref{T2} that problem $(\mathrm{P})$ admits nodal
solutions $(u_{1}^{\ast },u_{2}^{\ast })$ in $W_{0}^{1,p_{1}}(\Omega )\times
W_{0}^{1,p_{2}}(\Omega )$ where the nontrivial components $u_{1}^{\ast }$
and $u_{2}^{\ast }$ are not of the same constant sign. Assume that $%
u_{1}^{\ast }<0<u_{2}^{\ast }$. Test the first equation in $(\mathrm{P})$ by 
$-u_{1}^{\ast }{}^{-}$, in view of $(\mathrm{H}_{5})$, we get 
\begin{equation*}
\int_{\Omega }|\nabla u_{1}^{\ast -}|^{p_{1}}\text{ }\mathrm{d}%
x=-\int_{\Omega }f_{1}(x,u_{1}^{\ast },u_{2}^{\ast })u_{1}^{\ast }{}^{-}%
\text{ }\mathrm{d}x<0,
\end{equation*}%
which forces $u_{1}^{\ast -}=0$, a contradiction. So assume $u_{2}^{\ast
}<0<u_{1}^{\ast }$. Test the second equation in $(\mathrm{P})$ by $%
-u_{2}^{\ast }{}^{-}$, using $(\mathrm{H}_{5})$ it follows that 
\begin{equation*}
\int_{\Omega }|\nabla u_{2}^{\ast -}|^{p_{2}}\text{ }\mathrm{d}%
x=-\int_{\Omega }f_{2}(x,u_{1}^{\ast },u_{2}^{\ast })u_{2}^{\ast }{}^{-}%
\text{ }\mathrm{d}x<0.
\end{equation*}%
Hence, $u_{2}^{\ast -}=0$, a contradiction. Consequently, $u_{1}^{\ast }$
and $u_{2}^{\ast }$ cannot be of opposite constant sign. However,
considering Theorem \ref{T1} we can conclude that $u_{1}^{\ast }$ and $%
u_{2}^{\ast }$ must change sign simultaneously and therefore, (\ref{55}) is
satisfied. This completes the proof.
\end{proof}

\end{document}
